\subsection{TODO}

\paragraph{Immediate (blocking publication)}
None.

\paragraph{Near-term (needed for full benchmark)}
\begin{itemize}
    \item \textbf{Implement statistical analysis layer.} The benchmark runner produces raw per-scenario records (injection fired, consumption hits, propagation hits, recovery hits, downstream failure, task success) with $N$ repetitions per cell, but no statistical inference. We need confidence intervals, hypothesis testing (e.g., does one\_to\_many propagate failures significantly more than single\_origin?), effect sizes, and interaction analysis (topology $\times$ propagation mode $\times$ task family). Without this, we have counts but no way to draw causal conclusions or compare conditions rigorously.
    \item \textbf{Build and validate the graph construction module.} Convert enriched \texttt{Trace} objects (with \texttt{depends\_on} populated) into event-as-node graphs for GNN training, encoding node features (event type, agent role, tool name, metadata) and \texttt{depends\_on} edges as directed causal dependencies. Produce both static graph snapshots (per-event) and temporal event-stream formats.
    \item \textbf{Expand task evaluator coverage.} \texttt{CompetitiveIntelligenceEvaluator} does not exist at all — needs a deterministic rule-based evaluator matching the depth of \texttt{FinancialEvaluator} and \texttt{CodeReviewEvaluator}. \texttt{ResearchSynthesisEvaluator} exists but is a lightweight 5-check heuristic; needs expansion to proper depth. Without these, the benchmark runner cannot score task success for those families. IN-PROGRESS

    \item \textbf{Generate large-scale trace dataset.} Run the full cross-product matrix (all topologies $\times$ all task families $\times$ all LEP types $\times$ all propagation modes $\times$ repetitions) to produce a statistically powered dataset for GNN training and evaluation.

    \item \textbf{Diagrammatic representation of topologies for the paper}

\end{itemize}

\paragraph{Medium-term (experimental)}
\begin{itemize}
    \item \textbf{Generate large-scale trace dataset.} Current \texttt{pilot\_output/} has 24 executions and \texttt{out/} has 72. Target is a dataset covering all topologies ($\times$9), all task families ($\times$4), all LEP types ($\times$5), with benign, single-LEP, and counterfactual conditions, with repetition.
    \item \textbf{Train and evaluate GNN detectors.} Use the \texttt{ObservableExporter} traces (detector-visible only) to train GNNs for LEP detection, then evaluate on held-out traces. Compare against baselines (random, degree-based, temporal heuristics).
    \item \textbf{Add propagation-magnitude metrics.} Inspired by QUIVER: per-edge sensitivity classification (amplifier/absorber/threshold-sensitive), path-product approximation for transitive LEP amplification, and bifurcation thresholds (minimum perturbation strength causing structural graph changes).
\end{itemize}

\subsection{Details}

\paragraph{System overview.}
AgentGraph-v3 is a benchmark for studying how Local Execution Perturbations (LEPs) cause downstream failures in multi-agent LLM systems and how those failures propagate across agent boundaries. A \texttt{ScenarioSpec} defines a task family, workspace fixture, execution topology, and a set of LEP configurations. The \texttt{ScenarioRunner} executes the scenario by running agents through a topology of stages, recording every event as a \texttt{TraceEvent} in a \texttt{Trace} (schema version 3.0.0). Paired traces share an \texttt{execution\_id} and differ only in \texttt{variant} (benign=\texttt{a}, malignant=\texttt{b}), enabling direct comparison. The trace schema enforces a strict \texttt{observable}/\texttt{hidden} split: the \texttt{ObservableExporter} strips all LEP-related keys and benchmark metadata for safe distribution to detector models, while the \texttt{AnalysisExporter} retains full labels for offline analysis.

\paragraph{Trace schema (v3.0.0).}
Each trace contains 11 event types: \texttt{user\_input}, \texttt{system\_init}, \texttt{topology\_transition}, \texttt{reasoning}, \texttt{tool\_call}, \texttt{tool\_result}, \texttt{agent\_handoff}, \texttt{memory\_retrieval}, \texttt{memory\_write}, \texttt{final\_response}, \texttt{llm\_output}. Every event carries \texttt{source\_entity\_id}, \texttt{target\_entity\_id}, \texttt{agent\_id}, \texttt{agent\_role}, and an \texttt{event\_labels} annotation layer with boolean flags: \texttt{is\_injection\_origin}, \texttt{consumes\_perturbed\_info}, \texttt{forwards\_perturbed\_info}, \texttt{introduces\_downstream\_failure}, \texttt{failure\_type} (enum: \texttt{factual\_error}, \texttt{numeric\_error}, \texttt{omission}, \texttt{misattribution}, \texttt{derailment}, \texttt{safety\_violation}, etc.). The \texttt{depends\_on} field stores causal dependency edges between events, populated at event emission time. Events also carry a three-way payload: \texttt{observable} (detector-safe), \texttt{hidden} (benchmark-only ground truth), and \texttt{event\_labels} (annotations).

\paragraph{LEP subsystem.}
Five LEP families are implemented in \texttt{leps/} and orchestrated by \texttt{LEPOrchestrator}:
\begin{itemize}
    \item \textbf{ToolResultCorruptionLEP}: Replaces the return value of a targeted \texttt{read\_text\_file} call with corrupted content. Seven canonical variants per task family (e.g., \texttt{numeric\_corruption} for financial analysis, \texttt{partial\_omission} for code review, \texttt{source\_swap} for research synthesis). Fires at the \texttt{tool\_result} boundary.
    \item \textbf{IndirectPromptInjectionLEP}: Embeds adversarial instructions inside retrieved file content (six variants: \texttt{ignore\_previous}, \texttt{exfiltrate}, \texttt{skip\_verification}, \texttt{change\_destination}, \texttt{trust\_this}, \texttt{terminate\_early}). Fires at the \texttt{tool\_result} boundary.
    \item \textbf{HandoffCorruptionLEP}: Mutates the \texttt{HandoffPayload} (summary, findings, attribution) in transit between agents. Six variants: \texttt{omit\_key\_finding}, \texttt{alter\_severity}, \texttt{replace\_value}, \texttt{remove\_uncertainty}, \texttt{unsupported\_recommendation}, \texttt{swap\_attribution}. Fires at the \texttt{agent\_handoff} boundary.
    \item \textbf{InputDisregardLEP}: Injects a behavioral instruction into the receiving agent's context, causing it to ignore or underweight the handoff. Five variants: \texttt{start\_scratch}, \texttt{discard\_warnings}, \texttt{ignore\_verifier}, \texttt{recency\_bias}, \texttt{prefer\_memory}. Fires at the \texttt{agent\_handoff} boundary.
    \item \textbf{MemoryPoisoningLEP}: Seeds false entries into the memory store for later retrieval by agents. Currently incompatible — no memory subsystem integration. Three pre-defined poisoned values per task family.
\end{itemize}
LEPs use boundary-aware routing: \texttt{BOUNDARY\_LEPS} maps \texttt{tool\_result} to \texttt{LEP\_TOOL\_RESULT\_CORRUPTION} and \texttt{LEP\_INDIRECT\_PROMPT\_INJECTION}; \texttt{agent\_handoff} to \texttt{LEP\_HANDOFF\_CORRUPTION} and \texttt{LEP\_INPUT\_DISREGARD}. One-shot semantics prevent re-firing after a successful mutation. The \texttt{TriggerMatcher} evaluates \texttt{InjectionTrigger} conditions (event type, source agent, tool name, path pattern, content pattern, occurrence count, event index bounds, probability) with idempotency enforcement.

\paragraph{Execution topologies.}
Three topologies are implemented, each designed to isolate a distinct
structural dimension of perturbation propagation. \texttt{review\_loop}
is a 2-stage sequential topology (researcher $\to$ analyst) with a bounded
backedge: the analyst can hand back to the researcher for revision, and the
review cycle is capped at one iteration. This topology isolates the effect
of a producer--critic feedback loop: the same perturbation can be observed
both at first consumption (researcher $\to$ analyst) and at revision
(analyst $\to$ researcher). \texttt{branch\_and\_verify} is a 3-stage
diamond topology where two independent branches (researcher, analyst) feed a
single verifier that reconciles their outputs. This topology isolates
redundancy and verification: the verifier has two independent sources and
can detect, reject, or recover from a perturbation that affects one branch.
\texttt{coordinator\_workers} is a 4-stage fan-in topology where a
coordinator decomposes a task and delegates to three specialist workers
(security, correctness, test), each of which hands back to the coordinator
for aggregation. This topology isolates many-to-one aggregation trust: a
single perturbed specialist can contaminate the coordinator's synthesis of
all three outputs.

\paragraph{Task families and evaluators.}
Three task families with deterministic rule-based evaluators:
\begin{itemize}
    \item \textbf{code\_review}: Two fixtures (\texttt{easy}, \texttt{conflicting}) with required security issues (path traversal, sanitize stub, validate access bypass, off-by-one). \texttt{CodeReviewEvaluator} keyword-matches output against required issues and checks for false-positive acceptance phrases.
    \item \textbf{financial\_analysis}: Two fixtures (\texttt{clean}, \texttt{version\_conflict}) with 9 required numeric facts (Q3 revenue \$1.52M, operating costs \$790K, etc.) and version conflict detection (v1 vs. v2). \texttt{FinancialEvaluator} extracts figures via regex and flags downstream failure when v1 values appear in v2-required output.
    \item \textbf{research\_synthesis}: One fixture (\texttt{conflicting}) with Paper A and Paper B. \texttt{ResearchSynthesisEvaluator} uses a 5-check weighted heuristic (retrieval, no fabrication, synthesis, report quality, gap identification) with pass threshold 0.5.
\end{itemize}

\paragraph{Exporters.}
Three export formats serve distinct audiences: \texttt{ObservableExporter} (detector-safe, strips \texttt{LEP\_*} keys, \texttt{hidden} fields, and \texttt{event\_labels}); \texttt{AnalysisExporter} (full trace with labels, LEP instances, propagation paths); \texttt{PrefixExporter} (event prefixes up to injection/failure milestones with binary labels for early-warning causal forecasting).

\paragraph{Current pilot results.}
A 24-execution pilot with \texttt{linear\_2} topology: \texttt{LEP\_TOOL\_RESULT\_CORRUPTION} is the only LEP demonstrating the full injection$\to$consumption$\to$propagation$\to$failure chain. \texttt{LEP\_HANDOFF\_CORRUPTION} and \texttt{LEP\_INPUT\_DISREGARD} fire but show zero consumption.

\subsection{Related Work}

\paragraph{Single-agent tool-use evaluation.}
The capability of LLM agents to select, sequence, and execute tools has been
benchmarked extensively. \textbf{AgentBench}~\citep{liu2023agentbench}
evaluates LLM-as-agent performance across eight interactive environments
(database, web, operating system, knowledge graph, game, embodied, and
others), establishing a standardized paradigm for measuring agent reasoning
and decision-making under tool access. \textbf{ToolBench}~\citep{qin2023toolbench}
and \textbf{API-Bank}~\citep{li2023apibank} extend this axis by constructing
tool-rich environments from real-world API suites, measuring whether agents
can retrieve, plan, and invoke the correct tools to fulfill complex
instructions. These benchmarks establish that tool-use reliability is a
first-class evaluation dimension, but they are inherently single-agent:
no agent-to-agent communication, handoffs, or shared state. XXX
inherits the tool-use evaluation paradigm and instruments individual tool
calls as potential LEP injection boundaries.

\paragraph{Multi-agent pipeline frameworks.}
Several frameworks formalize role-based multi-agent workflows with structured
inter-agent communication. \textbf{MetaGPT}~\citep{hong2023metagpt} encodes
standard operating procedures into an assembly-line pipeline (product
manager $\to$ architect $\to$ engineer $\to$ QA), using structured handoffs
and output verification between roles to prevent cascading hallucinations.
\textbf{AgentVerse}~\citep{chen2023agentverse} dynamically assembles
specialized expert agents (coder, reviewer, tester) for collaborative
task-solving. Both systems demonstrate that well-structured handoff contracts
improve multi-agent output quality, but neither systematically injects
failures at handoff boundaries or measures how those failures propagate
downstream. XXX treats inter-agent handoffs as typed experimental boundaries
at which controlled LEPs may be introduced.

\paragraph{Failure attribution and multi-agent tracing.}
Closest to XXX in spirit is \textbf{AgenTracer}~\citep{XXXX}, which
studies failure attribution in multi-agent trajectories. AgenTracer creates
failure trajectories using programmed fault injection and counterfactual
replay, then trains a model to attribute failures to particular agents or
steps. The research question it addresses is: \emph{which agent or step
caused a failure?} XXX asks a different question: given a controlled
local perturbation, \emph{how does its effect propagate through the
execution dependency graph}---is the perturbation consumed, propagated,
recovered from, or does it induce downstream failure? Where AgenTracer
focuses on localizing responsibility for observed failures, XXX
characterizes the propagation process between a controlled intervention
and its eventual outcome, including recovery.

\paragraph{Multi-agent execution graphs.}
\textbf{ForestBench}~\citep{XXXX} maps heterogeneous multi-agent execution
traces into unified collaboration graphs, enabling structural rather than
outcome-only evaluation of multi-agent systems. ForestBench also includes
perturbation studies, making it a close structural antecedent. However, its
objective is collaboration-quality evaluation: given an execution trace,
what does the collaboration graph reveal about team dynamics? XXX shares
the commitment to graph-based trace representation but uses the graph for
a different purpose---measuring perturbation propagation across typed
execution boundaries (tool results, handoffs, memory writes) rather than
evaluating collaboration structure.

\paragraph{Adversarial perturbation and propagation.}
In the single-agent setting, \textbf{ARE} (Agent Robustness Evaluation)
~\citep{ma2024are} models an agent as a computation graph of components
(vision encoder, policy LM, self-evaluator, tree-search value function)
and measures how adversarial perturbations propagate through that graph.
ARE finds that compromising the reflexion evaluator raises attack success
by~15\% and compromising the tree-search value function adds a further~20\%,
demonstrating that inference-time compute improving benign performance can
simultaneously introduce new vulnerabilities. ARE is a conceptual
precedent for modeling perturbation flow through a system graph, but its
graph is the internal component structure of a single agent operating
within one context window; it does not address handoff-level propagation
across distinct agents with separate reasoning traces, memory stores, and
tool interfaces. \textbf{AgentHarm}~\citep{andriushchenko2024agentharm}
evaluates whether jailbroken agents maintain harmful multi-step tool-use
capability, but operates in a single-agent setting without inter-agent
communication.

\paragraph{The gap.}
Existing work does not jointly provide the combination XXX targets:
typed interventions at runtime boundaries, explicit execution-dependency
graphs, and measurement of whether perturbations are consumed, propagated,
recovered from, or induce downstream failure across multi-agent topologies.
AgenTracer attributes failures post-hoc; ARE measures propagation within
single-agent component graphs; MetaGPT and AgentVerse provide multi-agent
pipeline architecture without perturbation instrumentation; ForestBench
enables structural trace evaluation without controlled intervention. XXX
fills this gap by providing a systematic protocol for studying how locally
injected failures propagate through directed multi-agent topologies,
with reproducible perturbation injection, typed boundary routing, and
paired controlled executions enabling attribution of downstream differences
to localized interventions.

\paragraph{LEP grounding papers.}
Each Local Execution Perturbation (LEP) implemented in this benchmark is
grounded in peer-reviewed or workshop papers that directly study or
demonstrate the corresponding attack vector in LLM or multi-agent systems.
Table~\ref{tab:lep_papers} summarizes the mapping; we elaborate on the
direct relevance of each paper below.

\subsubsection{LEP-to-paper grounding}

\paragraph{Tool Result Corruption.}
\textbf{InjecAgent}~\citep{liu2024injecagent} is a benchmark specifically
designed to evaluate indirect prompt injection attacks on LLM-powered agent
systems. It introduces a dynamic environment where agents must process tool
outputs that may contain embedded adversarial instructions, and measures
whether agents consume corrupted results or detect the injection. This
directly models our \texttt{LEP\_TOOL\_RESULT\_CORRUPTION}: a tool call
succeeds normally, but the returned content is silently altered (numeric
corruption, partial omission, source swap), causing the agent to act on
false information. InjecAgent's focus on tool-output trust boundaries makes
it the closest structural antecedent to our tool result corruption variants.

\paragraph{Indirect Prompt Injection.}
\textbf{Greshake et al.``Not What You've Signed Up For''}
~\citep{greshake2023not} is the foundational paper that defined
indirect prompt injection (IPI) as a distinct attack class. The authors
demonstrate that adversaries can remotely compromise LLM-integrated
applications by embedding malicious instructions in data sources the model
retrieves---web pages, documents, emails, code repositories---without ever
directly interacting with the LLM. Their taxonomy of attack impacts (data
exfiltration, information ecosystem contamination, worming, function
hijacking) directly underpins our six IIP variants
(\texttt{ignore\_previous}, \texttt{exfiltrate}, \texttt{skip\_verification},
\texttt{change\_destination}, \texttt{trust\_this}, \texttt{terminate\_early}).
The paper's demonstration that retrieved prompts act as ``arbitrary code
execution'' for the LLM is the core mechanism our IIP LEP exploits.

\textbf{Rahman \& Kim ``The Framing Gap''}~\citep{rahman2026framing}
extends Greshake et al.'s framework specifically to tool-using agents,
showing that IIP exfiltration attacks defeat surface-level defenses (input
sanitization, output filtering) that work in single-turn chatbot settings.
This directly motivates our focus on tool-using agent topologies: the
defenses that protect chatbots do not protect agents that retrieve and act
on documents via tools. The paper bridges the gap between foundational IIP
theory and the tool-using agent setting our benchmark targets.

\paragraph{Memory Poisoning.}
\textbf{ToolHazard}~\citep{mou2026toolhazard} is a security benchmark for
LLM-based agents that includes memory and context poisoning scenarios. It
evaluates whether agents can identify and disregard poisoned entries in
their memory store when those entries are retrieved and presented as trusted
prior knowledge. This directly models our
\texttt{LEP\_MEMORY\_POISONING}: false information is seeded into the memory
store and later retrieved by an agent that treats it as ground truth.
ToolHazard's inclusion of memory-specific attack scenarios makes it the
primary grounding for this LEP family.

\paragraph{Handoff Corruption.}
\textbf{MAgIC}~\citep{zhou2023magic} is the only dedicated benchmark for
multi-agent generative communication. It evaluates how well multi-agent
systems develop and use communication protocols for collaborative tasks,
explicitly testing information transfer between agents. The benchmark
documents three failure modes directly corresponding to our handoff
corruption variants: message loss (\texttt{omit\_key\_finding}), message
distortion (\texttt{alter\_severity}, \texttt{replace\_value}), and receiver
failure to act on messages (\texttt{remove\_uncertainty},
\texttt{unsupported\_recommendation}, \texttt{swap\_attribution}). MAgIC's
focus on inter-agent message fidelity provides the only peer-reviewed
empirical evidence that handoff corruption is a measurable, reproducible
phenomenon in multi-agent LLM systems.

\paragraph{Input Disregard.}
\textbf{HackAPrompt}~\citep{schulhoff2023hackaprompt} is a large-scale
benchmark of 14,134 adversarial prompts targeting LLM instruction-following.
Its core finding---that LLMs exhibit inconsistent ``instruction hierarchy''
and can be systematically induced to ignore, discard, or override their own
prior instructions and context---directly models all five Input Disregard
variants. \texttt{start\_scratch} maps to ``ignore everything above''
patterns; \texttt{discard\_warnings} maps to risk-flag suppression;
\texttt{ignore\_verifier} maps to correction override; \texttt{recency\_bias}
maps to context window manipulation favoring the latest input; and
\texttt{prefer\_memory} maps to models trusting their own prior context over
received input. HackAPrompt demonstrates that input disregard is not a
corner case but a systematic vulnerability across model families.

\subsection{Task-family benchmarks}

\paragraph{code\_review.}
We ground the \texttt{code\_review} task family in
\textbf{SWE-bench}~\citep{jimenez2023swebench}, a large-scale benchmark of
real-world GitHub issues requiring code comprehension, reasoning about
correctness, and producing patches. SWE-bench is the standard reference for
evaluating LLM-based software engineering agents and provides the canonical
paradigm for agentic code reasoning tasks.

\paragraph{financial\_analysis.}
We ground the \texttt{financial\_analysis} task family in
\textbf{ConvFinQA}~\citep{chen2022convfinqa}, a conversational financial QA
benchmark requiring multi-step numerical reasoning over tables. ConvFinQA
directly evaluates chain-of-numerical-reasoning capabilities that
multi-agent financial analysis systems must exhibit when extracting and
comparing financial figures across document versions.

\paragraph{research\_synthesis.}
We ground the \texttt{research\_synthesis} task family in
\textbf{GAIA}~\citep{mialon2023gaia}, a general AI assistant benchmark that
evaluates whether agents can autonomously gather, synthesize, and act on
information from multiple sources. GAIA's information-synthesis dimension
maps directly onto the research-synthesis workflow: agents retrieve
conflicting documents, reconcile findings, and produce a unified report.

\subsubsection{Topology precedents}

\paragraph{review\_loop.}
The \texttt{review\_loop} topology implements the producer--critic--revise
pattern described by \textbf{Reflexion}~\citep{shinn2023reflexion}, where
one actor produces an output, a separate evaluator reviews it, and
linguistic feedback conditions a revision on the next turn. Reflexion
demonstrates that a bounded review cycle (actor $\to$ critic $\to$ revise)
improves task success rates across code and reasoning tasks, providing the
canonical precedent for our 2-agent backedge topology.

\paragraph{branch\_and\_verify.}
The \texttt{branch\_and\_verify} topology follows the independent-branch
verification pattern of \textbf{MAVEN}~\citep{yao2026maven}, a
multi-agent verification--elaboration network where independent reasoning
branches produce outputs that a central judge reconciles. MAVEN explicitly
separates logical defense from factual grounding across specialized agent
roles, motivating our design where two independent branches (researcher,
analyst) feed a verifier that detects and resolves disagreement.

\paragraph{coordinator\_workers.}
The \texttt{coordinator\_workers} topology implements the
coordinator--delegate--aggregate pattern of \textbf{MetaGPT}
~\citep{hong2023metagpt}, where a central coordinator decomposes tasks,
fans them out to specialist workers, and aggregates results through
structured verification gates. MetaGPT's assembly-line paradigm with
explicit role separation and inter-stage validation is the closest
architectural precedent for our topology, in which a coordinator fans out
to multiple specialists and reconciles their handoffs.

\subsection{Methodology}

\Yue{We may need a clear setup here, including: What platform are we using (like your reproducibility and infra below, with GPU specified)? What tasks we are doing (we should build our experiments on some real multi-agent task benchmarks---datasets)? How do we set up the agents? What matrices are we evaluating? You may refer to the example below.}

\color{teal}
\textbf{Execution Platform and Infrastructure.}
We implement our multi-agent workflows using \textit{[orchestration framework]}, with LLM inference served through \textit{[serving framework]}.
Our experiments use \textit{[LLM model(s)]} with \textit{[precision / quantization / context length]}.
All experiments are conducted on \textit{[GPU model and number of GPUs]}, with \textit{[CPU / memory / other relevant hardware information]}.
Unless otherwise stated, all agents in a workflow use the same LLM backend to isolate the effects of workflow structure and abnormality propagation. We report the full software and hardware configuration to support reproducibility.

\textbf{Multi-Agent Tasks and Workloads.}
We evaluate \OurMethod{} using both real-world multi-agent benchmarks and
controlled synthesized workloads. For real workloads, we use \textit{[Benchmark A]}, \textit{[Benchmark B]}, and \textit{[Benchmark C]}, covering tasks such as \textit{[tool use / retrieval / planning / coding / coordination]}.
These benchmarks provide realistic task objectives, tool interactions, and standardized success criteria.

To complement the real workloads, we also construct synthesized multi-agent tasks that allow us to systematically control workflow structure and risk propagation patterns. Each synthesized task is drawn from one of four task families---\texttt{code\_review}, \texttt{financial\_analysis}, \texttt{research\_synthesis}, \texttt{competitive\_intelligence}---paired with workspace fixtures, and instantiated through a \texttt{ScenarioSpec} that combines a task family, fixture ID, execution topology, and LEP configuration. Each instantiation produces a set of agents with assigned roles (e.g., inspector, reviewer, analyst, synthesizer, searcher), tool access (\texttt{list\_directory}, \texttt{read\_text\_file}, \texttt{write\_file}, \texttt{search\_files}, \texttt{handoff}, \texttt{submit\_final}, \texttt{read\_memory}, \texttt{write\_memory}), a memory configuration, and a target task outcome grounded in fixture-specific ground truth. We vary the number of agents (2--5, depending on topology), workflow topology, communication pattern, information-sharing policy, memory configuration, and task dependency depth (2--5 stages) to produce workflows with different propagation characteristics.

The injected abnormalities correspond to five LEP families---\texttt{LEP\_TOOL\_RESULT\_CORRUPTION}, \texttt{LEP\_INDIRECT\_PROMPT\_INJECTION}, \texttt{LEP\_HANDOFF\_CORRUPTION}, \texttt{LEP\_INPUT\_DISREGARD}, and \texttt{LEP\_MEMORY\_POISONING}---and are configured to produce isolated, one-to-many, many-to-one, or multi-hop propagation patterns depending on the injection boundary (\texttt{tool\_result}, \texttt{agent\_handoff}, or \texttt{memory\_write}) and the downstream topology structure.
This enables precise ground-truth annotation of abnormality sources, propagation paths, and downstream outcomes, and allows us to distinguish true risk accumulation from simple co-occurrence.

\color{black}


\paragraph{Experimental design.}
Three topologies serve as the structural substrate for propagation analysis.
\texttt{review\_loop} (2-stage with bounded revision cycle) isolates the
effect of a producer--critic feedback loop; \texttt{branch\_and\_verify}
(3-stage diamond) isolates redundancy and verification; and
\texttt{coordinator\_workers} (4-stage fan-in) isolates many-to-one
aggregation trust. Each topology targets a distinct structural dimension of
perturbation propagation: revision-cycle sensitivity in review loops,
detection-and-recovery via redundant branches in branch-and-verify, and
aggregation trust in coordinator-workers. Each scenario is defined by a
\texttt{ScenarioSpec} with four parameters: \texttt{task\_family}
(code\_review, financial\_analysis, research\_synthesis), \texttt{fixture\_id}
(workspace fixture), \texttt{topology} (one of three), and \texttt{condition}
(\texttt{benign}, \texttt{single\_lep}, \texttt{counterfactual}). For each
task--fixture pair, we generate one benign trace, one single-LEP trace per
LEP type, and one counterfactual trace (same as single-LEP but with LEPs
removed). Paired traces share an \texttt{execution\_id} and differ only in
\texttt{variant}, enabling within-scenario causal attribution: if benign
passes and malignant fails on the same task and topology, the LEP is
causally implicated.

\paragraph{LEP injection.}
Canonical operators are selected deterministically per (task\_family, LEP\_code) pair via \texttt{get\_canonical\_operator()}, ensuring reproducible injection across runs. The \texttt{LEPOrchestrator} registers LEP instances from \texttt{LEPConfig} objects and evaluates them at boundary events using \texttt{evaluate\_for\_boundary()}, which routes to eligible LEPs based on event type. Triggers use \texttt{TriggerMatcher} for semantic condition matching with idempotency (a trigger fires at most once per scenario) and occurrence counting.

\paragraph{LLM backend.}
The primary execution backend is \texttt{APIBackend}, which makes HTTP requests to the OpusMax API gateway (\texttt{LLM\_BASE\_URL}) using only the Python stdlib \texttt{urllib} — no Anthropic SDK required. It supports native \texttt{tool\_use} mode with system prompt, conversation history, and tool definitions. The default model is \texttt{claude-sonnet-5} with temperature 0.1. Model selection is configurable via \texttt{LLM\_MODEL}.

\paragraph{Trace annotation and evaluation.}
After execution, each trace is annotated with \texttt{EventLabels} recording the perturbation lifecycle. Task-specific evaluators (\texttt{CodeReviewEvaluator}, \texttt{FinancialEvaluator}, \texttt{ResearchSynthesisEvaluator}) check the final output against fixture ground truth to determine \texttt{task\_success} and \texttt{downstream\_failure}.

\paragraph{Reproducibility and infrastructure.}
LLM inference is served through two separate backends: (1) the OpusMax API gateway for \texttt{claude-sonnet-5}, contacted via stdlib \texttt{urllib}; and (2) a \texttt{vllm} server hosting \texttt{/} (4-bit quantized), deployed on the DeltaAI cluster (GH200 120GB GPU). All experiments are run on the DeltaAI cluster login node \texttt{gh-login01.delta.ncsa.illinois.edu} (kernel 6.4.0-150600.23.53-64kb, x86\_64) with 144 CPUs and 250 GiB RAM. LLM inference for the open-source model path is served through \texttt{vllm} hosting \texttt{/} (4-bit quantized) on GH200 120GB GPUs. All agents in a workflow use the same backend to isolate the effects of workflow structure and abnormality propagation. The codebase uses \texttt{torch} for the graph encoding module (encoder, exporter, graph builder). Seeds are set per scenario via \texttt{ScenarioBuildConfig.seed}. Traces are serialized as JSON files (schema v3.0.0) and also exported as JSONL via the three exporters. Pilot results and experimental outputs are stored in \texttt{pilot\_output/} and \texttt{out/} respectively.

\paragraph{Graph construction (planned).}
Once \texttt{depends\_on} is populated internally, each trace is converted into an event-as-node graph: every execution event becomes a node, and \texttt{depends\_on} causal dependencies become directed edges (A$\to$B when B depends on A). This preserves fan-in structure and enables path-based propagation analysis directly on the execution graph. Node features carry event type, agent role, tool name, stage-local index, and observable metadata; LEP labels (injection origin, consumption, propagation, failure) are node-level annotations for training. Graph-level labels indicate whether the execution ultimately experienced a downstream failure.

\bibliography{references}
\nocite{*}
